% Source: https://artofproblemsolving.com/community/c3h2325619_try_this_on_for_practice
% Stable identity: aops-topic-2325619-q01
\textbf{(a)} Since $1000=1001-1\equiv-1\pmod7$, a number is congruent modulo $7$ to the alternating sum of its three-digit groups, starting from the right with a $+$. The $78$ digits form $26$ groups; from the right they are
\[336,\,494,\,161,\,518,\,141,\,163,\,605,\,382,\,904,\,074,\,279,\,564,\,837,\,852,\,907,\,687,\,008,\,985,\,570,\,423,\,195,\,316,\,237,\,089,\,792,\,115.\]
Taking them in pairs (first minus second):
\[-158-357-22+223+830-285-15+220-977+147-121+148+677=310.\]
Since $310=44\cdot7+2$, the remainder is $\boxed{2}$.

\textbf{(b)} $N\equiv2\pmod7$, so $N^N\equiv2^N\pmod7$. Since $2^3=8\equiv1\pmod7$, only $N\bmod3$ matters. The digit sum of $N$ is $348$, a multiple of $3$, so $3\mid N$ and $2^N=\left(2^3\right)^{\dfrac N3}\equiv1\pmod7$. The remainder is $\boxed{1}$.
